\section{Balance de la empresa: el punto más bajo, el auge y la siguiente etapa}

Esta sección convierte la línea emocional de la entrevista en un balance de la empresa emergente. El «qué satisfacción» de Chen Mian (陈冕) viene de una victoria parcial, pero la siguiente etapa plantea preguntas más difíciles: cómo se organiza el equipo, cómo se estabiliza el flujo de caja, cómo se retiene a los usuarios del producto, cómo se sostiene el modelo de negocio y cómo se gana un lugar en la mente profesional de los diseñadores. Las empresas emergentes de aplicaciones de IA suelen enfrentar un segundo examen después del auge: cuando llegan los usuarios, ¿el sistema soporta la carga? Cuando llegan los ingresos, ¿la organización puede crecer? Cuando llegan los competidores, ¿siguen en pie las barreras de entrada?

\begin{knowledgebox}{El segundo examen después del auge}
\begin{tabularx}{\textwidth}{p{0.20\textwidth}p{0.34\textwidth}X}
\toprule
Examen & Pregunta clave & Señal de aprobación \\
\midrule
Examen de producto & ¿Los usuarios usan una y otra vez el mismo flujo de trabajo? & Retención, número de proyectos, profundidad de la colaboración. \\
Examen de negocio & ¿Existe una razón estable para pagar? & Suscripciones, contratos de equipo, conversión a pago. \\
Examen de organización & ¿El equipo puede pasar del modo de combate a un funcionamiento sistemático? & Ritmo de investigación y desarrollo, atención al cliente, entrega, disciplina financiera. \\
Examen de competencia & ¿Qué queda cuando los grandes modelos y los competidores te alcanzan? & Barreras de datos, marca, activos y ecosistema. \\
\bottomrule
\end{tabularx}
\end{knowledgebox}

\begin{warningbox}{Después del auge, no persigas solo más visibilidad}
Cuando una aplicación de IA se vuelve viral, es fácil que el equipo siga persiguiendo la difusión, pero el verdadero peligro es que las capacidades básicas no estén a la altura: si la estabilidad, el costo, los derechos de autor, la colaboración y el servicio posventa no acompañan el crecimiento, la visibilidad terminará amplificando los problemas del producto.
\end{warningbox}

\subsection{Cómo convertir las emociones del fundador en capacidad organizacional}

Un CEO combativo puede atravesar el punto más bajo, pero una empresa no puede depender por mucho tiempo solo de la energía emocional de su fundador. La siguiente etapa exige convertir el criterio del fundador en mecanismos organizacionales: hoja de ruta del producto, sistema de retroalimentación de usuarios, disciplina financiera, evaluación de modelos, procesos de derechos de autor y éxito del cliente. De lo contrario, un auge puede quedarse en un pico pasajero.

\begin{knowledgebox}{De la energía personal a los mecanismos organizacionales}
\begin{tabularx}{\textwidth}{p{0.22\textwidth}p{0.34\textwidth}X}
\toprule
Energía personal & Ya organizada & Función \\
\midrule
Intuición de producto & Investigación de usuarios y sistema de evaluación & Evita depender solo de la decisión del fundador. \\
Modo de combate & Gestión del ritmo y mecanismos de priorización & Evita que el equipo se agote a largo plazo. \\
Presión de financiamiento & Modelo financiero y gestión del runway & Evita volver a quedarse sin efectivo. \\
Entusiasmo de los usuarios & Éxito del cliente y gestión de la comunidad & Convierte la emoción en retención. \\
\bottomrule
\end{tabularx}
\end{knowledgebox}

\subsection{Resumen de la sección}

La siguiente etapa de Lovart no consiste en demostrar que «puede volverse popular», sino en demostrar que «puede consolidarse como mesa de trabajo de diseño». Esto exige pasar de las emociones del fundador a los mecanismos organizacionales, y del Magic Moment a una entrega repetible.
